\section{Related Work}

% Plan:
% - introductory work on HAR -> initial paper that introduces our dataset
% - outline work done on HAR -> deep neural networks / etc (not related to our project
% - introduce the concept of missing data 
% - outline work done on missing data -> hossain 2018/20, xaviar, reyes-ortiz, saha
% - outline general sensor faults
% - mention our purpose again

HAR using wearable inertial sensors is a research area with well-established benchmark datasets
and models. Lara and Labrador \cite{lara-2013} provided a survey on the research area, covering
sensors used, feature extraction methods for time series based datasets, learning models typically
used, and evaluation metrics. The UCI HAR dataset \cite{reyes-ortiz-2013} introduced by Reyes-Ortiz
et al. \cite{reyes-ortiz-2013} provided a multiclass SVM baseline achieving 96\% accuracy. Deep
learning methods used by Ronao et al. \cite{ronao-2016} provided similar results with 94.79\%
classification accuracy on the raw sensor data, and 95.75\% classification accuracy if the featured
frequency information is included. Saha et al. \cite{saha-2025} provided another baseline using
a lightweight LSTM model with KNN imputation, achieving, 93.89\% accuracy. These results confirm
that HAR on clean data is a well-studied problem.

% missing imputation discussion below
Missing sensor readings is the most common studied problem when regarding non-clean data
assumptions in sensor produced data. Adhikari et al. \cite{adhikari-2022} provided a
survey on imputation as a tecnique for replacing missing data in internet of things systems.
In the context of HAR, Hossain et al. \cite{hossain-2018} provided a study on sensor-based HAR with
missing data by evaluating the effect of randomly missing accelerometer data, finding that accuracy
rates decreased significantly even with small missing data rates. Hossain et al. \cite{hossain-2020}
later introduced the idea of explicitly inducing missing data in the training sets to make models
more resistant to missing data and showed that recognition accuracy significantly increased without
requiring imputation. Saha et al. \cite{saha-2025} addressed the same problem on the UCI HAR
dataset, using KNN imputation to recover missing readings. Xaviar et al. \cite{xaviar-2023}
proposed a cleaning module: a multimodal fusion model with a denoising autoencoder that learns to
correct missing input signals before classification.

These works represent important contributions in sensor degraded data, however they suffer from the
same limitation of only considering sensor dropout, and thus study imputation as a solution to this
problem. However many sensor fault types cannot be simply solved by imputation as they do not
represent missing data.

Sensor corruption literature has been studied extensively. Ni et al. \cite{ni-2009} provided
a taxonomy of sensor network data faults drawn from real deployments: stuck-at-faults, offset
faults, calibration faults, noise, drift, etc. Li et al. \cite{li-2020} extended this in a review
of sensor fault diagnosis techniques, distinguishing between incipient failures that are difficult
to detect early (bias, drift, gain) and abrupt failures such as sudden noise spikes or hardware
disconnection. The discussion of sensor degraded data on HAR datasets have largely been focused on
abrupt failures, for which data imputation can fill the gaps, however incipient failures ...

We base our sensor corruption framework on the physical sensor faults introduced by Li et al.
\cite{li-2020}. We focus on sensor bias, drift, gain (amplification, attenuation), noise, and
dropout which correspond to a documented physical failure mode. Thus we ensure our framework
and corruption scenarios we study are not arbitrary but reflect degradation patterns that real HAR
deployments will encounter.

% citations cut:
%     - samara-2008: sensor faults but in aircraft systems... didn't seem too applicable but
%         mentioned in li-2020
%     - qi-2021: impacts of dirty data in general... could be used to argue that our investigation
%         is important but I couldn't find anything concrete
%
% there should only be 10 citations then

% all citations just to test -> there should be 12 total citations at the end
% \cite{samara-2008} \cite{ni-2009} \cite{reyes-ortiz-2013} \cite{lara-2013} \cite{ronao-2016}
% \cite{hossain-2018} \cite{hossain-2020} \cite{li-2020} \cite{qi-2021} \cite{adhikari-2022}
% \cite{xaviar-2023} \cite{saha-2025}
